\section{Discussion}
\label{sec:discussion}

\subsection{What transfers across domains, and what does not}

Three patterns held in both domains. First, difficulty follows phase
dynamics. Short, manoeuvre-dominated phases (takeoff; departure and arrival)
are harder than long, steady ones (cruise; transit), which supports
phase-specific memories in both adapters. Second, grounding matters more
than the generator. In both domains, retrieval-only baselines capture part
of the gain, and a guarded LLM on top of retrieval captures the rest. A
maritime-only check reinforces the phase point directly: a two-parameter
dead-reckoning baseline that only holds the last trusted speed and reads
the live trusted heading already matches or beats TrajRAG and AIS-ACNet in
Houston departure/transit and Los Angeles/Long Beach transit, and is
markedly worse than both in the turning-dominated arrival phase
(Section~\ref{sec:simple-baselines}); the genuinely hard, differentiating
comparisons across every method tested are concentrated in manoeuvre-heavy
phases, not steady ones. Third,
memories can be extended across routes and corridors without loss of
accuracy. One thing did \emph{not} transfer automatically: the numeric value
of the feasibility tolerance. This distinction between procedure and
constant is the main practical lesson for anyone porting the framework to a
further domain, such as rail, road freight or unmanned aircraft.

\subsection{When is a trained domain model preferable?}

No significant difference was detected between AIS-ACNet and TrajRAG in any
of the three Houston phases (Section~\ref{sec:significance}), though
AIS-ACNet had the lower sample mean in all three; it is also cheaper per
inference. We read matching, rather than beating, a strong ship-specific
network as the relevant result: AIS-ACNet is purpose-trained per phase on
this corridor's own traffic, while TrajRAG uses no maritime training data
at all. When a single domain and a single port are of interest, when a
curated training set and a machine-learning pipeline are available, and when
the traffic pattern is stable, a dedicated network remains a sound
engineering choice, and on this evidence it is not dominated by TrajRAG in
accuracy. TrajRAG is preferable when one service must cover several domains
or many corridors, when new corridors must come online quickly, when
training data per corridor are limited, or when operators need an
interpretable trail of the retrieved precedents and of the guard's
interventions -- advantages of engineering flexibility, not of raw
accuracy, on the evidence reported here.

\subsection{Operational integration}

In deployment, TrajRAG sits behind an interference detector and in front of
the operator display (Fig.~\ref{fig:architecture}). The reconstructed track
should be presented as an estimate, with the guard's low-confidence flag
shown whenever a proposal was clamped. When trusted reports resume, or an
independent source such as surface radar or multilateration becomes
available, the reconstruction should be replaced by that source rather than
fused with it. For VTS centres that already operate coastal radar, the
reconstruction is most valuable outside radar coverage, or as a cross-check
that flags a broadcast track that departs from where precedent and physics
say the ship should be.

\subsection{Threats to validity}
\label{sec:threats}

\textbf{Internal validity.} (i) The maritime reference-speed rule now holds
the reference speed at the last trusted SOG for the whole blackout
(Section~\ref{sec:maritime-speedrule}), which reads no post-onset field; an
earlier variant that updated it from the recorded SOG at each reconstructed
step is reported only as a brief, superseded pre-correction comparison
(Section~\ref{sec:r1-result}), not as a result. (ii) The two domains' main
results use different LLM backbones (gpt-5.1 and gpt-4.1) and slightly
different retrieval sizes; an earlier, pre-correction pass
(Section~\ref{sec:maritime-backbone}) measured the effect directly and
found it phase-dependent rather than uniform, but that pass was not
re-run under the preprocessing and retrieval-memory corrections, so it is
disclosed as a directional check rather than a fully confirmed
confound-free result. (iii) The modern forecasters were evaluated on
a shorter horizon and a subset of voyages; Table~\ref{tab:item4-matched}
gives TrajRAG and AIS-ACNet a matched comparison on that same subset and
horizon, but the forecasters themselves were not retrained under the
corrected data. (x) Budget constraints limited TrajRAG's headline
evaluation to 50 Houston and 10 Los Angeles/Long Beach voyages rather than
the full 125-per-corridor test set. Every other method is rescored on the
identical subset for a matched comparison, and AIS-ACNet's full-125 and
matched-subset means agree closely, but the Los Angeles/Long Beach cell
size in particular ($n=10$) gives limited statistical power, and several
corridor comparisons that might reach significance at $n=125$ do not at
$n=10$. \textbf{External validity.} (iv)
Both maritime corridors are in the United States, and only cargo and tanker
vessels were included. (v) The trusted-signal assumption fails on vessels
whose heading comes from a satellite compass. Such vessels would need a
different trusted signal, or none, in which case the guard alone bounds the
error. (vi) The blackouts are simulated on real trajectories. Unlike
aviation, where multilateration offered an independent reference for real
interference episodes, we do not yet have an independently verified
maritime incident. \textbf{Construct validity.} (vii) MHE weights all steps
equally. Operators may care more about end-of-blackout error or about
cross-track error in a channel. Normalised errors
(Table~\ref{tab:normalized}) partly address the first concern. (viii) An
earlier implementation of TrajRAG's retrieval memory had two problems: it
excluded only the exact query voyage from the corridor corpus, not the
other 124 test voyages, so a reconstruction could retrieve another test
voyage as a precedent; and it had no chronological cutoff among the
remaining (non-test) voyages, so a reconstruction could retrieve a
precedent recorded after the query. We treat both as correctness bugs
rather than disclosable design choices and fixed them directly
(Section~\ref{sec:data}): every result in this paper is reconstructed
against a memory built only from the chronologically earliest half of each
corridor's non-test voyages, with every test voyage additionally excluded,
so no retrieved precedent can postdate or belong to the test set, by
construction. (ix) The maritime phase label used throughout evaluation
comes from offline, full-voyage SOG segmentation
(Section~\ref{sec:maritime-adapter}), not from the online,
trusted-information-only phase rule $\pi$ that
Algorithm~\ref{alg:reconstruction} specifies; online phase identification
for maritime is unimplemented and untested. (xi) Every family of related
comparisons in this paper (the Houston main set, the Los Angeles/Long Beach
main set, the short-horizon matched comparison, and the ablations) is now
corrected for multiple testing with Holm's method within that family
(Section~\ref{sec:significance}); this changes two substantive conclusions
relative to treating each raw $p$-value in isolation. First, several
Los Angeles/Long Beach comparisons that look significant uncorrected
(every Kalman cell, departure ARIMA) do not survive correction at $n=10$
(Table~\ref{tab:wilcoxon-la}). Second, the maritime ablation's
retrieval-removal effect, which looked significant in transit and arrival
before correction, does not survive it at $n=20$
(Table~\ref{tab:maritime-ablation}); we report its direction as suggestive
only. In neither case does a non-significant corrected result establish
that the compared methods, or the ablated and full pipelines, perform
equivalently -- only that these sample sizes do not resolve a difference
once corrected for testing several related hypotheses together.
